\documentclass{article}
\PassOptionsToPackage{numbers,sort,compress}{natbib}
\usepackage[main,final]{neurips_2025}
\usepackage{fontspec}
\usepackage[russian]{babel}
\setmainfont{cmunrm.otf}[BoldFont=cmunbx.otf,ItalicFont=cmunti.otf,BoldItalicFont=cmunbi.otf]
\setsansfont{cmunss.otf}
\setmonofont{cmuntt.otf}
\usepackage[hidelinks]{hyperref}
\usepackage{url}
\usepackage{booktabs}
\usepackage{amsmath,amssymb}
\makeatletter
\renewcommand{\@noticestring}{Исследовательский проект.}
\makeatother
\renewenvironment{abstract}{\vskip 0.075in\centerline{\large\bfseries\abstractname}\vspace{0.5ex}\begin{quote}}{\par\end{quote}\vskip 1ex}
\setlength{\emergencystretch}{2em}
\title{Адаптивное условное диффузионное сэмплирование при поступлении новых наблюдений в процессе генерации}
\author{}
\begin{document}
\maketitle
\begin{abstract}
Данная работа посвящена исследованию способов адаптации диффузионных моделей для генерации траекторий к условиям, когда новый контекст поступает во время генерации. В задачах прогнозирования движения дополнительные наблюдения могут уточнить возможные варианты будущего или существенно изменить их вероятности, однако учёт этих данных может потребовать повторных вычислений, которые не всегда допустимы при ограниченном времени и вычислительном бюджете. Я исследую, как обновлять формируемую траекторию с учётом новых данных, в каких случаях можно использовать уже выполненные вычисления и как выбранный способ обновления влияет на качество прогноза и затраты на его получение. Для этого я рассмотрю существующие подходы к диффузионному прогнозированию, потоковой генерации и условному сэмплированию, сопоставлю их предположения и ограничения и выделю способы, применимые к изменению контекста внутри незавершённого процесса генерации. Сравнение предполагается проводить при разных моментах поступления наблюдений, разной степени их влияния на прогноз и сопоставимых вычислительных бюджетах, учитывая точность траекторий, разнообразие возможных исходов и время получения обновлённого результата. При этом я отдельно рассмотрю качество вероятностного прогноза, поскольку близость одной сгенерированной траектории к наблюдаемому движению ещё не означает, что модель корректно отражает неопределённость будущего. На основе выявленных ограничений я предложу возможные способы улучшения существующих подходов и проверю, в каких условиях они позволяют эффективнее использовать новые наблюдения. Исследование направлено на определение возможностей и границ адаптации диффузионного сэмплирования к изменяющемуся контексту при ограниченных вычислительных ресурсах.
\end{abstract}
\noindent\textbf{Ключевые слова:} диффузионные модели; генерация траекторий; вероятностное прогнозирование; условное сэмплирование; потоковые наблюдения; вычислительная эффективность.

\clearpage
\section{Обзор литературы}
\subsection{Диффузионная генерация и прогнозирование траекторий}

Диффузионные модели формируют результат через последовательное расшумление. Базовая дискретная формулировка представлена в DDPM \citep{p37}, а непрерывная связь между стохастическим процессом и дифференциальным уравнением — в Score-SDE \citep{p38}. Для моей темы важно, что между начальным шумом и готовой траекторией существует промежуточное состояние. При поступлении новых данных возникает вопрос, можно ли использовать это состояние для обновлённого прогноза и какие ошибки при этом появляются.

В прогнозировании движения MID использует историю траекторий и взаимодействия участников для условной генерации \citep{p09}. LED дополнительно обучает инициализацию, позволяющую пропускать часть расшумления \citep{p10}, а MotionDiffuser моделирует совместные варианты движения нескольких участников и поддерживает направленное сэмплирование \citep{p11}. Эти работы дают предметную основу для сравнения качества, разнообразия и стоимости прогнозов. DiffTraj также работает с траекториями, но решает задачу генерации синтетических GPS-данных \citep{p12}. Поэтому я рассматриваю её как источник идей о представлении данных, а не как прямой аналог обновления прогноза.

Отдельную группу составляют методы планирования. Diffuser и Decision Diffuser используют условную генерацию для построения поведения с заданными целями и ограничениями \citep{p14}, \citep{p15}. Diffusion Policy генерирует последовательности действий по наблюдениям \citep{p17}; Diffusion Planner рассматривает автономное движение \citep{p13}, а DiffPhyCon — управление физическими системами \citep{p18}. Здесь качество определяется не только соответствием распределению данных, но и результатом исполнения выбранных действий. AdaptDiffuser дополнительно изменяет сам планировщик через дообучение на отобранных синтетических данных \citep{p16}. Я отделяю такую адаптацию модели от обновления одного незавершённого прогноза.

\subsection{Новые наблюдения и повторное использование вычислений}

Наиболее близки к теме работы, сохраняющие часть генерации между поступлениями данных. Streaming Diffusion Policy поддерживает буфер действий с разными уровнями шума и уточняет его по новым наблюдениям \citep{p01}. CL-DiffPhyCon также использует асинхронное расшумление, включая обратную связь при генерации дальнейших управляющих воздействий \citep{p02}. Они показывают, что использование нового контекста до завершения всей последовательности уже является самостоятельным направлением исследований. Diffusion Forcing даёт более общий механизм обучения последовательностей с разными уровнями шума у элементов \citep{p05}.

RTI-DP использует ранее предсказанную последовательность как начальное приближение для следующего расчёта \citep{p04}. Sequential Flow Matching рассматривает перенос предыдущего распределения в обновлённое по наблюдениям \citep{p03}. В отличие от повторного использования отдельной траектории, здесь обучается переход между распределениями. Авторы последней работы отмечают накопление ошибок при рекурсивных обновлениях и обсуждают адаптивное повторное зашумление и перезапуск как дальнейшие направления. Это конкретная связь с моей темой, но не основание считать такую общую идею новой.

Задержку можно учитывать и на уровне исполнения. RTC согласует соседние фрагменты действий при асинхронной генерации \citep{p06}. Reactive Diffusion Policy выделяет быстрый контур обратной связи \citep{p07}, а FA-RDP адаптирует частоту вывода и число шагов генерации к неоднозначности действий \citep{p08}. Эти подходы показывают, что реакция на новые данные может достигаться разными механизмами. Поэтому я буду различать обновление промежуточного сэмпла, перенос готового прогноза и коррекцию исполнения: близкие практические цели не делают эти задачи одинаковыми.

\subsection{Условное распределение по наблюдениям}

Для анализа обновлений полезны работы с временными рядами. CSDI и SSSD восстанавливают пропущенные значения по доступным наблюдениям \citep{p19}, \citep{p23}. TimeGrad строит авторегрессионный прогноз \citep{p20}, а TimeDiff генерирует будущий участок совместно и использует специальные механизмы задания условия \citep{p21}. TSDiff вводит условия при выводе и рассматривает уточнение готовых прогнозов \citep{p24}. Эти методы помогают выбрать основу для экспериментов, но применение к наблюдениям, поступающим внутри расшумления, требует собственного протокола.

Score-based Data Assimilation отделяет модель наблюдений от обучения генератора и восстанавливает траектории по неполным или шумным измерениям \citep{p22}. В обратных задачах похожую роль играют DDRM и Diffusion Posterior Sampling \citep{p29}, \citep{p28}. Однако согласование с измерениями может опираться на приближения, а ограничения линейной модели наблюдений подходят не для любого контекста. Кроме того, Classifier-Free Guidance управляет соответствием условию и разнообразием \citep{p40}; усиление такого соответствия само по себе не подтверждает корректность вероятностного обновления.

\subsection{Вычислительная эффективность и повторное зашумление}

Сокращение затрат не обязательно связано с адаптацией контекста. DDIM, DPM-Solver, DPM-Solver++ и UniPC ускоряют численное сэмплирование \citep{p30}, \citep{p31}, \citep{p32}, \citep{p33}. EDM систематизирует выбор параметризации, расписания шума и сэмплера \citep{p36}. Progressive Distillation и Consistency Models уменьшают число шагов через изменение обучения \citep{p35}, \citep{p34}, а Flow Matching представляет смежный способ обучать перенос распределений \citep{p39}. Эти работы нужны, чтобы отделить возможную пользу обновления контекста от обычного ускорения генератора и учитывать стоимость дополнительного обучения.

Повторное зашумление также является известным инструментом. RePaint использует повторное сэмплирование при восстановлении изображений \citep{p27}, Restart исследует сочетание добавления шума и обратного интегрирования \citep{p25}, а RePS переносит restart-подход на обратные задачи \citep{p26}. Их результаты обосновывают включение подобных операций в сравнение, но не устанавливают, какая коррекция нужна после конкретного нового наблюдения в задаче траекторий.

\subsection{Направление исследования}

Рассмотренные работы уже предлагают способы потоковой генерации, повторного использования прогнозов и ускорения сэмплирования. Поэтому я не формулирую задачу как создание первого метода, учитывающего новые наблюдения. Я исследую, как эти способы работают при изменении контекста внутри незавершённой генерации траектории и как их применимость зависит от момента поступления данных, величины изменения прогноза и доступного бюджета. Для сравнения важны и точность, и сохранение неопределённости, и время получения обновления. На основе выявленных ограничений я предложу возможные улучшения; конкретный подход и его отличие от ближайших методов будут определены по результатам дальнейшего анализа и экспериментов.


\clearpage
\bibliographystyle{unsrtnat}
\bibliography{references}
\end{document}
